\chapter{Проектирование тестов} \label{ch2}

В данной главе описаны техники проектирования тестов, применённые в работе.
Для каждой техники сформированы проверяемые условия (test cover items), контрольные примеры (test cases) и оптимизированные наборы тестов.
Дополнительно приведено сравнение выбранных методов по покрытию, трудоёмкости и практической ценности.

\section{Проектирование тестов на основе выбранных технологий} \label{ch2:sec1}

Стандарт ISO/IEC/IEEE 29119-4:2021 выделяет три группы техник проектирования тестов: методы на основе спецификации, методы на основе структуры тестируемого объекта и методы на основе опыта тестировщика~\cite{iso29119-4}.
В данной работе сознательно используются первые две группы.
Методы на основе спецификации применяются к требованиям, форматам входных данных и ожидаемым результатам, то есть не требуют знания исходного кода.
Методы на основе структуры применяются к исходному коду и позволяют проверить операторы, решения, ветви и другие элементы управляющего потока.

Для проектирования тестов приложения выбраны следующие техники на основе спецификации:
\begin{itemize}
    \item разбиение на классы эквивалентности (Equivalence Partitioning, EP);
    \item анализ граничных значений (Boundary Value Analysis, BVA).
\end{itemize}

Из техник на основе структуры выбраны:
\begin{itemize}
    \item тестирование операторов (Statement Testing);
    \item тестирование ветвей (Branch Testing).
\end{itemize}

Таким образом, EP и BVA используются как техники класса specification-based, соответствующие группе техник из приложения B стандарта, а Statement Testing и Branch Testing используются как техники класса structure-based, соответствующие группе техник из приложения C.
Такое разделение важно для трассируемости: входные и выходные классы проектируются по требованиям к приложению, а внутренние условия алгоритма Манакера проверяются структурными тестами по коду.

\subsection{Метод разбиения на классы эквивалентности} \label{ch2:sec1-subsec1}

Метод EP в ISO/IEC/IEEE 29119-4:2021 относится к техникам на основе спецификации и предполагает разбиение входов или выходов тестируемого объекта на классы, которые должны обрабатываться одинаково.
В данной работе EP применяется к пользовательским входам, JSON-вводу и наблюдаемым результатам анализа строки.
Для алгоритма Манакера такими классами являются строки разной длины и различные семантические типы строк: полностью палиндромные, непалиндромные, содержащие чётные или нечётные максимальные палиндромы, строки с повторяющимися символами и т.д.
Внутренние состояния алгоритма, например нахождение позиции внутри интервала $[l,r]$, не рассматриваются как самостоятельные классы эквивалентности; они используются далее в структурных техниках Statement Testing и Branch Testing.

Проверяемые условия для метода EP приведены в \taref{tab:ep-test-cover-items}.

\noindent
\begingroup
\centering
\small
\begin{longtable}[c]{|p{2.4cm}|p{8.8cm}|p{2.3cm}|p{1.7cm}|}
    \caption{Проверяемые условия с использованием метода EP}
    \label{tab:ep-test-cover-items}
    \\
    \hline
    Имя проверки & Содержание проверки & Слой данных & Результат \\ \hline
    \endfirsthead
    \captionsetup{format=tablenocaption,labelformat=continued}
    \caption[]{}\\
    \hline
    Имя проверки & Содержание проверки & Слой данных & Результат \\ \hline
    \endhead
    \hline
    \endfoot
    \hline
    \endlastfoot
    EP-TCI-1 & Строка пустая & Входные & 0 \\ \hline
    EP-TCI-2 & Строка длины 1 & Входные & 1 \\ \hline
    EP-TCI-3 & Строка длины 2 без чётного палиндрома & Входные & 1 \\ \hline
    EP-TCI-4 & Строка длины 2 с чётным палиндромом & Входные & 1 \\ \hline
    EP-TCI-5 & Строка без палиндромов длины $>1$ & Входные & 1 \\ \hline
    EP-TCI-6 & Строка, целиком являющаяся нечётным палиндромом & Входные & 1 \\ \hline
    EP-TCI-7 & Строка, целиком являющаяся чётным палиндромом & Входные & 1 \\ \hline
    EP-TCI-8 & Строка с несколькими перекрывающимися нечётными палиндромами & Входные & 1 \\ \hline
    EP-TCI-9 & Строка с несколькими перекрывающимися чётными палиндромами & Входные & 1 \\ \hline
    EP-TCI-10 & Строка из одинаковых символов & Входные & 1 \\ \hline
    EP-TCI-11 & Строка с пробелами и специальными символами & Входные & 1 \\ \hline
    EP-TCI-12 & Корректный JSON-файл содержит поле \texttt{input} & Входные & 1 \\ \hline
    EP-TCI-13 & JSON-файл не содержит обязательного поля \texttt{input} & Входные & 0 \\ \hline
    EP-TCI-14 & Файл по пути отсутствует & Входные & 0 \\ \hline
    EP-TCI-15 & Наибольший палиндром имеет нечётную длину & Выходные & 1 \\ \hline
    EP-TCI-16 & Наибольший палиндром имеет чётную длину & Выходные & 1 \\ \hline
    EP-TCI-17 & Все чётные радиусы равны 0 & Выходные & 1 \\ \hline
    EP-TCI-18 & Присутствуют ненулевые чётные радиусы & Выходные & 1 \\ \hline
    EP-TCI-19 & Формируется рекомендация о полной симметрии строки & Выходные & 1 \\ \hline
    EP-TCI-20 & Формируется рекомендация о низкой палиндромной плотности & Выходные & 1 \\ \hline
    EP-TCI-21 & Результаты двух алгоритмов совпадают & Выходные & 1 \\ \hline
    EP-TCI-22 & Результаты двух алгоритмов расходятся & Выходные & 0 \\ \hline
    EP-TCI-23 & Результат успешно сохраняется в JSON-файл & Выходные & 1 \\ \hline
\end{longtable}
\normalsize
\endgroup

Контрольные примеры для метода EP приведены в \taref{tab:ep-test-cases}.

\noindent
\begingroup
\centering
\small
\begin{longtable}[c]{|p{2cm}|p{7.2cm}|p{6.3cm}|}
    \caption{Контрольные примеры для метода EP}
    \label{tab:ep-test-cases}
    \\
    \hline
    Test case & Входные данные & Ожидаемый вывод \\ \hline
    \endfirsthead
    \captionsetup{format=tablenocaption,labelformat=continued}
    \caption[]{}\\
    \hline
    Test case & Входные данные & Ожидаемый вывод \\ \hline
    \endhead
    \hline
    \endfoot
    \hline
    \endlastfoot
    EP-TC-1 & \texttt{""} & Выброс исключения или сообщение о недопустимом вводе \\ \hline
    EP-TC-2 & \texttt{"a"} & $d_{odd}=[1]$, $d_{even}=[0]$, наибольший палиндром \texttt{"a"} \\ \hline
    EP-TC-3 & \texttt{"ab"} & $d_{odd}=[1,1]$, $d_{even}=[0,0]$ \\ \hline
    EP-TC-4 & \texttt{"aa"} & наибольший палиндром \texttt{"aa"}, чётный радиус в центре ненулевой \\ \hline
    EP-TC-5 & \texttt{"abc"} & все радиусы минимальны, рекомендация о слабой симметрии \\ \hline
    EP-TC-6 & \texttt{"abacaba"} & наибольший палиндром совпадает со всей строкой \\ \hline
    EP-TC-7 & \texttt{"abba"} & наибольший палиндром чётной длины 4 \\ \hline
    EP-TC-8 & \texttt{"ababa"} & несколько перекрывающихся нечётных палиндромов \\ \hline
    EP-TC-9 & \texttt{"abbaabba"} & несколько крупных чётных палиндромов \\ \hline
    EP-TC-10 & \texttt{"aaaaa"} & возрастающие радиусы, максимальное число палиндромов \\ \hline
    EP-TC-11 & \texttt{"a-b-a"} & специальные символы обрабатываются как обычные \\ \hline
    EP-TC-12 & JSON с \texttt{"input":"level"} & корректное чтение строки из файла \\ \hline
    EP-TC-13 & JSON без поля \texttt{input} & ошибка валидации формата \\ \hline
    EP-TC-14 & несуществующий путь к файлу & ошибка чтения файла \\ \hline
\end{longtable}
\normalsize
\endgroup

В результате оптимизации 14 базовых тест-кейсов сокращены до 8 оптимизированных при сохранении покрытия существенных входных и выходных классов эквивалентности.
Оптимизированный набор представлен в \taref{tab:ep-test-cases-optimized}.

\noindent
\begingroup
\centering
\small
\begin{longtable}[c]{|p{1.6cm}|p{6.2cm}|p{3cm}|p{4.8cm}|}
    \caption{Оптимизированные контрольные примеры для метода EP}
    \label{tab:ep-test-cases-optimized}
    \\
    \hline
    Test case & Входные данные & Покрываемые EP-TCI & Ожидаемый вывод \\ \hline
    \endfirsthead
    \captionsetup{format=tablenocaption,labelformat=continued}
    \caption[]{}\\
    \hline
    Test case & Входные данные & Покрываемые EP-TCI & Ожидаемый вывод \\ \hline
    \endhead
    \hline
    \endfoot
    \hline
    \endlastfoot
    EP-TC-1 & \texttt{""} & 1 & Ошибка ввода \\ \hline
    EP-TC-2 & \texttt{"a"} & 2, 15 & минимальный корректный результат \\ \hline
    EP-TC-3 & \texttt{"aa"} & 4, 16, 18 & чётный палиндром длины 2 \\ \hline
    EP-TC-4 & \texttt{"abc"} & 3, 5, 17, 20 & только одиночные палиндромы \\ \hline
    EP-TC-5 & \texttt{"abacaba"} & 6, 8, 19, 21 & полная симметрия строки \\ \hline
    EP-TC-6 & \texttt{"abbaabba"} & 7, 9, 16, 18 & крупные чётные палиндромы \\ \hline
    EP-TC-7 & \texttt{"aaaaa"} & 10 & максимальная палиндромная плотность \\ \hline
    EP-TC-8 & JSON / ошибка файла & 11--14, 22, 23 & проверка ввода-вывода \\ \hline
\end{longtable}
\normalsize
\endgroup

\subsection{Метод анализа граничных значений} \label{ch2:sec1-subsec2}

Метод BVA в ISO/IEC/IEEE 29119-4:2021 является техникой на основе спецификации и ориентирован на границы классов эквивалентности.
В работе он применяется прежде всего к границам входных и выходных доменов: минимальной допустимой длине строки, строке длины 2, пустому JSON-файлу, минимальной строке в JSON и крайним значениям длины найденного палиндрома.
Границы индексов в алгоритме Манакера дополнительно используются как обоснование структурных тестов, поскольку они связаны с риском ошибок типа off-by-one в реализации.

Проверяемые условия для метода BVA приведены в \taref{tab:bva-test-cover-items}.

\noindent
\begingroup
\centering
\small
\begin{longtable}[c]{|p{2.4cm}|p{8.8cm}|p{2.3cm}|p{1.7cm}|}
    \caption{Проверяемые условия с использованием метода BVA}
    \label{tab:bva-test-cover-items}
    \\
    \hline
    Имя проверки & Содержание проверки & Тип границы & Результат \\ \hline
    \endfirsthead
    \captionsetup{format=tablenocaption,labelformat=continued}
    \caption[]{}\\
    \hline
    Имя проверки & Содержание проверки & Тип границы & Результат \\ \hline
    \endhead
    \hline
    \endfoot
    \hline
    \endlastfoot
    BVA-TCI-1 & Длина строки $0$ & Входные & 0 \\ \hline
    BVA-TCI-2 & Длина строки $1$ & Входные & 1 \\ \hline
    BVA-TCI-3 & Длина строки $2$ без совпадения символов & Входные & 1 \\ \hline
    BVA-TCI-4 & Длина строки $2$ с совпадением символов & Входные & 1 \\ \hline
    BVA-TCI-5 & Минимальный нечётный палиндром длины 3 & Входные & 1 \\ \hline
    BVA-TCI-6 & Палиндром касается левой границы строки & Граница строки & 1 \\ \hline
    BVA-TCI-7 & Палиндром касается правой границы строки & Граница строки & 1 \\ \hline
    BVA-TCI-8 & Центр палиндрома находится в начале массива проверки & Алгоритмическая & 1 \\ \hline
    BVA-TCI-9 & Центр палиндрома находится в конце массива проверки & Алгоритмическая & 1 \\ \hline
    BVA-TCI-10 & Максимальный радиус обновляет правую границу ровно до $n-1$ & Алгоритмическая & 1 \\ \hline
    BVA-TCI-11 & Проверка условия выхода за левую границу даёт ложь на шаге $-1$ & Алгоритмическая & 1 \\ \hline
    BVA-TCI-12 & Проверка условия выхода за правую границу даёт ложь на шаге $n$ & Алгоритмическая & 1 \\ \hline
    BVA-TCI-13 & Пустой JSON-файл & Входные & 0 \\ \hline
    BVA-TCI-14 & JSON-файл с минимально допустимой строкой & Входные & 1 \\ \hline
    BVA-TCI-15 & Путь сохранения результата указывает на новый файл & Выходные & 1 \\ \hline
    BVA-TCI-16 & Наибольший палиндром длины 1 & Выходные & 1 \\ \hline
    BVA-TCI-17 & Наибольший палиндром длины, равной длине всей строки & Выходные & 1 \\ \hline
\end{longtable}
\normalsize
\endgroup

Контрольные примеры BVA приведены в \taref{tab:bva-test-cases}.

\noindent
\begingroup
\centering
\small
\begin{longtable}[c]{|p{2cm}|p{7cm}|p{6.5cm}|}
    \caption{Контрольные примеры для метода BVA}
    \label{tab:bva-test-cases}
    \\
    \hline
    Test case & Входные данные & Ожидаемый вывод \\ \hline
    \endfirsthead
    \captionsetup{format=tablenocaption,labelformat=continued}
    \caption[]{}\\
    \hline
    Test case & Входные данные & Ожидаемый вывод \\ \hline
    \endhead
    \hline
    \endfoot
    \hline
    \endlastfoot
    BVA-TC-1 & \texttt{""} & ошибка ввода \\ \hline
    BVA-TC-2 & \texttt{"a"} & минимально допустимый результат \\ \hline
    BVA-TC-3 & \texttt{"ab"} & два одиночных палиндрома \\ \hline
    BVA-TC-4 & \texttt{"aa"} & минимальный чётный палиндром \\ \hline
    BVA-TC-5 & \texttt{"aba"} & минимальный нечётный палиндром длины 3 \\ \hline
    BVA-TC-6 & \texttt{"baaab"} & палиндром касается обеих границ строки \\ \hline
    BVA-TC-7 & \texttt{"abccba"} & правая граница обновляется до конца строки \\ \hline
    BVA-TC-8 & пустой JSON-файл & ошибка чтения \\ \hline
    BVA-TC-9 & JSON с \texttt{"input":"x"} & корректная загрузка минимальной строки \\ \hline
    BVA-TC-10 & сохранение в новый файл & успешная запись результата \\ \hline
\end{longtable}
\normalsize
\endgroup

Оптимизированный набор BVA-тестов приведён в \taref{tab:bva-test-cases-optimized}.

\noindent
\begingroup
\centering
\small
\begin{longtable}[c]{|p{1.8cm}|p{6.6cm}|p{2.8cm}|p{4.5cm}|}
    \caption{Оптимизированные контрольные примеры для метода BVA}
    \label{tab:bva-test-cases-optimized}
    \\
    \hline
    Test case & Входные данные & Покрываемые BVA-TCI & Ожидаемый вывод \\ \hline
    \endfirsthead
    \captionsetup{format=tablenocaption,labelformat=continued}
    \caption[]{}\\
    \hline
    Test case & Входные данные & Покрываемые BVA-TCI & Ожидаемый вывод \\ \hline
    \endhead
    \hline
    \endfoot
    \hline
    \endlastfoot
    BVA-TC-1 & \texttt{""} & 1, 13 & ошибка ввода \\ \hline
    BVA-TC-2 & \texttt{"a"} & 2, 14, 16 & минимальный корректный сценарий \\ \hline
    BVA-TC-3 & \texttt{"aa"} & 3, 4 & проверка длины 2 \\ \hline
    BVA-TC-4 & \texttt{"aba"} & 5, 8 & нечётный палиндром минимальной длины \\ \hline
    BVA-TC-5 & \texttt{"baaab"} & 6, 7, 10, 11, 12, 17 & палиндром на границах строки \\ \hline
    BVA-TC-6 & \texttt{"abccba"} & 9, 10 & чётный палиндром на правой границе \\ \hline
    BVA-TC-7 & сохранение в новый файл & 15 & успешная запись результата \\ \hline
\end{longtable}
\normalsize
\endgroup

\subsection{Тестирование операторов} \label{ch2:sec1-subsec3}

Метод Statement Testing относится к техникам на основе структуры и ориентирован на то, чтобы каждый значимый исполняемый оператор программы был выполнен хотя бы одним тестом.
Для приложения поиска палиндромов это включает:
\begin{itemize}
    \item инициализацию массивов радиусов;
    \item присваивание стартового радиуса;
    \item цикл расширения палиндрома;
    \item обновление границ $l,r$;
    \item суммирование радиусов при подсчёте количества палиндромов;
    \item формирование JSON-результата.
\end{itemize}
Перечисленные условия являются структурными test coverage items: они не подменяют классы эквивалентности, а дополняют их анализом исполняемых операторов.

Проверяемые условия ST приведены в \taref{tab:st-test-cover-items}.

\noindent
\begingroup
\centering
\small
\begin{longtable}[c]{|p{2.4cm}|p{10.2cm}|p{1.8cm}|}
    \caption{Проверяемые условия с использованием Statement Testing}
    \label{tab:st-test-cover-items}
    \\
    \hline
    Имя проверки & Содержание проверки & Результат \\ \hline
    \endfirsthead
    \captionsetup{format=tablenocaption,labelformat=continued}
    \caption[]{}\\
    \hline
    Имя проверки & Содержание проверки & Результат \\ \hline
    \endhead
    \hline
    \endfoot
    \hline
    \endlastfoot
    ST-TCI-1 & Инициализация массива $d_{odd}$ выполнена & 1 \\ \hline
    ST-TCI-2 & Инициализация массива $d_{even}$ выполнена & 1 \\ \hline
    ST-TCI-3 & Выполнено присваивание стартового радиуса вне интервала $[l,r]$ & 1 \\ \hline
    ST-TCI-4 & Выполнено присваивание стартового радиуса по зеркальной позиции & 1 \\ \hline
    ST-TCI-5 & Выполнен хотя бы один шаг расширения палиндрома & 1 \\ \hline
    ST-TCI-6 & Выполнена запись нового значения радиуса в массив & 1 \\ \hline
    ST-TCI-7 & Выполнено обновление границ $l,r$ & 1 \\ \hline
    ST-TCI-8 & Выполнен проход по всем позициям строки & 1 \\ \hline
    ST-TCI-9 & Выполнено восстановление нечётного максимального палиндрома & 1 \\ \hline
    ST-TCI-10 & Выполнено восстановление чётного максимального палиндрома & 1 \\ \hline
    ST-TCI-11 & Выполнено суммирование радиусов для подсчёта количества палиндромов & 1 \\ \hline
    ST-TCI-12 & Выполнено формирование рекомендации пользователю & 1 \\ \hline
\end{longtable}
\normalsize
\endgroup

Контрольные примеры ST приведены в \taref{tab:st-test-cases}.

\noindent
\begingroup
\centering
\small
\begin{longtable}[c]{|p{2cm}|p{7cm}|p{6.3cm}|}
    \caption{Контрольные примеры для Statement Testing}
    \label{tab:st-test-cases}
    \\
    \hline
    Test case & Входные данные & Ожидаемый вывод \\ \hline
    \endfirsthead
    \captionsetup{format=tablenocaption,labelformat=continued}
    \caption[]{}\\
    \hline
    Test case & Входные данные & Ожидаемый вывод \\ \hline
    \endhead
    \hline
    \endfoot
    \hline
    \endlastfoot
    ST-TC-1 & \texttt{"abacaba"} & выполняются основные операторы алгоритма Манакера \\ \hline
    ST-TC-2 & \texttt{"abba"} & выполняется ветка чётного палиндрома и восстановление результата \\ \hline
    ST-TC-3 & сохранение результата в JSON & выполняется запись файла и формирование отчёта \\ \hline
\end{longtable}
\normalsize
\endgroup

\subsection{Тестирование ветвей} \label{ch2:sec1-subsec4}

Метод Branch Testing относится к техникам на основе структуры и ориентирован на покрытие ключевых ветвей условий в управляющем потоке программы.
Для реализации алгоритма Манакера существенны следующие развилки:
\begin{itemize}
    \item позиция находится внутри интервала $[l,r]$ / вне интервала;
    \item расширение возможно / невозможно;
    \item границы обновляются / не обновляются;
    \item максимальный палиндром чётный / нечётный;
    \item входной JSON корректен / некорректен.
\end{itemize}
Эти условия относятся к структуре управляющего потока и поэтому рассматриваются в рамках техник приложения C, а не как отдельные классы эквивалентности.

Проверяемые условия BT приведены в \taref{tab:bt-test-cover-items}.

\noindent
\begingroup
\centering
\small
\begin{longtable}[c]{|p{2.4cm}|p{10.2cm}|p{1.8cm}|}
    \caption{Проверяемые условия с использованием Branch Testing}
    \label{tab:bt-test-cover-items}
    \\
    \hline
    Имя проверки & Содержание проверки & Результат \\ \hline
    \endfirsthead
    \captionsetup{format=tablenocaption,labelformat=continued}
    \caption[]{}\\
    \hline
    Имя проверки & Содержание проверки & Результат \\ \hline
    \endhead
    \hline
    \endfoot
    \hline
    \endlastfoot
    BT-TCI-1 & Ветвь обработки пустой строки & 0 \\ \hline
    BT-TCI-2 & Ветвь обработки корректной непустой строки & 1 \\ \hline
    BT-TCI-3 & Ветвь $i > r$ & 1 \\ \hline
    BT-TCI-4 & Ветвь $i \le r$ & 1 \\ \hline
    BT-TCI-5 & Ветвь успешного расширения палиндрома & 1 \\ \hline
    BT-TCI-6 & Ветвь без расширения & 1 \\ \hline
    BT-TCI-7 & Ветвь обновления границ & 1 \\ \hline
    BT-TCI-8 & Ветвь без обновления границ & 1 \\ \hline
    BT-TCI-9 & Ветвь выбора нечётного максимального палиндрома & 1 \\ \hline
    BT-TCI-10 & Ветвь выбора чётного максимального палиндрома & 1 \\ \hline
    BT-TCI-11 & Ветвь корректной загрузки JSON & 1 \\ \hline
    BT-TCI-12 & Ветвь ошибки формата JSON & 0 \\ \hline
\end{longtable}
\normalsize
\endgroup

Контрольные примеры BT приведены в \taref{tab:bt-test-cases}.

\noindent
\begingroup
\centering
\small
\begin{longtable}[c]{|p{2cm}|p{7cm}|p{6.3cm}|}
    \caption{Контрольные примеры для Branch Testing}
    \label{tab:bt-test-cases}
    \\
    \hline
    Test case & Входные данные & Ожидаемый вывод \\ \hline
    \endfirsthead
    \captionsetup{format=tablenocaption,labelformat=continued}
    \caption[]{}\\
    \hline
    Test case & Входные данные & Ожидаемый вывод \\ \hline
    \endhead
    \hline
    \endfoot
    \hline
    \endlastfoot
    BT-TC-1 & \texttt{""} & ошибка ввода \\ \hline
    BT-TC-2 & \texttt{"abc"} & покрытие ветви без расширения \\ \hline
    BT-TC-3 & \texttt{"aaaa"} & покрытие ветви с зеркальной позицией и расширением \\ \hline
    BT-TC-4 & \texttt{"abba"} & покрытие ветви выбора чётного максимального палиндрома \\ \hline
    BT-TC-5 & некорректный JSON & покрытие ветви ошибки формата \\ \hline
\end{longtable}
\normalsize
\endgroup

\section{Сравнительный анализ методов проектирования тестов} \label{ch2:sec2}

Преимущества и недостатки выбранных методов приведены в \taref{tab:methods-advantages}.

\noindent
\begingroup
\centering
\small
\begin{longtable}[c]{|p{3.4cm}|p{6.1cm}|p{6.1cm}|}
    \caption{Преимущества и недостатки методов проектирования тестов}
    \label{tab:methods-advantages}
    \\
    \hline
    Метод & Преимущества & Недостатки \\ \hline
    \endfirsthead
    \captionsetup{format=tablenocaption,labelformat=continued}
    \caption[]{}\\
    \hline
    Метод & Преимущества & Недостатки \\ \hline
    \endhead
    \hline
    \endfoot
    \hline
    \endlastfoot
    EP &
    Хорошо покрывает смысловые типы строк, сокращает число тестов за счёт выбора представителей классов, применим без готовой реализации &
    Может пропустить граничные дефекты, требует аккуратного выделения классов эквивалентности \\ \hline
    BVA &
    Эффективно выявляет ошибки на границах длины и индексов, хорошо дополняет EP &
    Быстро увеличивает число тестов при росте числа параметров, хуже работает без чётких диапазонов \\ \hline
    Statement Testing &
    Даёт минимальный структурный контроль, позволяет быстро проверить исполнение всех операторов &
    Не гарантирует покрытие альтернативных ветвей и исключительных путей \\ \hline
    Branch Testing &
    Проверяет ключевые логические развилки, хорошо дополняет спецификационные методы &
    Требует знания реализации, не покрывает все возможные пути исполнения \\ \hline
\end{longtable}
\normalsize
\endgroup

Сводное сравнение покрытий, количества тестов и трудозатратности приведено в \taref{tab:methods-comparison}.

\noindent
\begingroup
\centering
\small
\begin{longtable}[c]{|p{4.8cm}|p{2.3cm}|p{2.5cm}|p{2.5cm}|p{2.5cm}|}
    \caption{Сравнение методов проектирования тестов по ключевым метрикам}
    \label{tab:methods-comparison}
    \\
    \hline
    Метрика & EP & BVA & Statement Testing & Branch Testing \\ \hline
    \endfirsthead
    \captionsetup{format=tablenocaption,labelformat=continued}
    \caption[]{}\\
    \hline
    Метрика & EP & BVA & Statement Testing & Branch Testing \\ \hline
    \endhead
    \hline
    \endfoot
    \hline
    \endlastfoot
    Количество TCI & 30 & 17 & 12 & 12 \\ \hline
    Количество TC & 14 & 10 & 3 & 5 \\ \hline
    Количество оптимизированных TC & 8 & 7 & 2--3 & 4--5 \\ \hline
    Трудозатратность проектирования & Высокая & Средняя & Низкая & Средняя \\ \hline
    Требует знания реализации & Нет & Нет & Частично & Да \\ \hline
    Покрытие входных данных & Высокое & Высокое на границах & Низкое & Низкое \\ \hline
    Покрытие кодовых путей & Низкое & Низкое & Базовое & Хорошее \\ \hline
    Обнаружение граничных ошибок & Среднее & Высокое & Низкое & Среднее \\ \hline
\end{longtable}
\normalsize
\endgroup

EP обеспечивает наиболее широкий охват входных сценариев, включая разные типы строк, форматы ввода и виды результата.
BVA даёт особенно высокую ценность для алгоритма Манакера, так как многие дефекты сосредоточены на границах строки и при переходе между соседними центрами.
Statement Testing полезен как минимальный структурный контроль, а Branch Testing обеспечивает проверку важнейших логических развилок реализации.

\section{Проектирование тестов для загрузки строки и вывода результата} \label{ch2:sec3}

Отдельно спроектированы тесты для чтения строки из JSON-файла и для вывода результата в консоль и файл.
Для этих компонентов основным методом выступает EP с элементами BVA.

Проверяемые условия для чтения JSON-файла приведены в \taref{tab:io-tci}.

\noindent
\begingroup
\centering
\small
\begin{longtable}[c]{|p{3cm}|p{9.5cm}|p{2.8cm}|}
    \caption{Проверяемые условия для чтения JSON-файла и вывода результата}
    \label{tab:io-tci}
    \\
    \hline
    Имя проверки & Содержание проверки & Результат \\ \hline
    \endfirsthead
    \captionsetup{format=tablenocaption,labelformat=continued}
    \caption[]{}\\
    \hline
    Имя проверки & Содержание проверки & Результат \\ \hline
    \endhead
    \hline
    \endfoot
    \hline
    \endlastfoot
    IO-TCI-1 & JSON содержит поле \texttt{input} непустой строкой & 1 \\ \hline
    IO-TCI-2 & JSON существует, но поле \texttt{input} отсутствует & 0 \\ \hline
    IO-TCI-3 & Файл по указанному пути не существует & 0 \\ \hline
    IO-TCI-4 & Вывод в консоль содержит массивы радиусов и наибольший палиндром & 1 \\ \hline
    IO-TCI-5 & JSON-результат содержит все обязательные поля & 1 \\ \hline
    IO-TCI-6 & Рекомендация успешно записывается в выходной файл & 1 \\ \hline
\end{longtable}
\normalsize
\endgroup

Контрольные примеры для ввода-вывода приведены в \taref{tab:io-tc}.

\noindent
\begingroup
\centering
\small
\begin{longtable}[c]{|p{1.8cm}|p{8.1cm}|p{5.5cm}|}
    \caption{Контрольные примеры для ввода-вывода}
    \label{tab:io-tc}
    \\
    \hline
    Test case & Входные данные & Ожидаемый результат \\ \hline
    \endfirsthead
    \captionsetup{format=tablenocaption,labelformat=continued}
    \caption[]{}\\
    \hline
    Test case & Входные данные & Ожидаемый результат \\ \hline
    \endhead
    \hline
    \endfoot
    \endlastfoot
    IO-TC-1 & JSON: \texttt{\{"input":"level"\}} & строка корректно загружается \\ \hline
    IO-TC-2 & JSON без \texttt{input} & ошибка валидации \\ \hline
    IO-TC-3 & несуществующий файл & ошибка чтения файла \\ \hline
    IO-TC-4 & результат анализа \texttt{"abba"} & корректный вывод в консоль \\ \hline
    IO-TC-5 & сохранение результата \texttt{"abacaba"} & корректный JSON-файл результата \\ \hline
\end{longtable}
\normalsize
\endgroup

\section{Выводы} \label{ch2:conclusion}

В данной главе спроектированы тесты для консольного приложения, реализующего алгоритм Манакера и эталонный алгоритм расширения от центра.
Для проверки выбраны четыре техники проектирования тестов: EP, BVA, Statement Testing и Branch Testing.
Первые две техники относятся к группе specification-based и применяются к входам, выходам и наблюдаемому поведению приложения; вторые две техники относятся к группе structure-based и применяются к коду, операторам, условиям и ветвям алгоритма.
Сформированы проверяемые условия, контрольные примеры и оптимизированные наборы тестов, а также выполнен сравнительный анализ методов по покрытию и трудозатратности.
